\subsection{সুবিধা, খরচ ও সীমাবদ্ধতা}
\label{ch2:c10-12}
\begin{tabularx}{\linewidth}{|p{.28\linewidth}|X|X|}\hline
\textbf{বিষয়} & \textbf{প্রাসঙ্গিক সুবিধা} & \textbf{সীমা/খরচ}\\\hline
সম্পদ ও রক্ষণাবেক্ষণ & নিজের সব ভৌত সার্ভার কিনতে হয় না; অনেক অংশ সেবাদাতা পরিচালনা করে & সেবাদাতার ওপর নির্ভরশীলতা; মডেল অনুযায়ী নিয়ন্ত্রণ ও নিজস্ব দায়িত্ব বদলায়\\\hline
চাহিদা ও খরচ & দ্রুত সম্পদ বাড়ানো বা কমানো; প্রাথমিক বিনিয়োগ কমতে পারে & দীর্ঘদিন বেশি সংরক্ষণ, স্থানান্তর ও প্রক্রিয়াকরণে বিল বাড়ে\\\hline
প্রবেশ ও যৌথ কাজ & উপযুক্ত নেটওয়ার্কে বিভিন্ন ডিভাইস থেকে একই ডকুমেন্টে কাজ & নেটওয়ার্ক বা সেবাদাতার সেবা বন্ধ থাকলে কাজ ব্যাহত; সব কাজ অফলাইনে করা যায় না\\\hline
ডেটা ব্যবস্থাপনা & ভাগাভাগি ও পুনরুদ্ধারের ব্যবস্থা সহজ করা যায় & গোপনীয়তা, ভুল করে মুছে ফেলা ও অননুমোদিত ভাগাভাগির ঝুঁকি; আগের সংস্করণ বা ব্যাকআপ দরকার\\\hline
পরিবর্তন ও স্থানান্তর & নতুন সম্পদ দ্রুত নেওয়া যায় & সেবাদাতা বদলাতে ডেটা ও টুল সরানো জটিল হতে পারে; Vendor Lock-in\\\hline
\end{tabularx}
স্থানীয় ব্যবস্থায় প্রতিষ্ঠান হার্ডওয়্যার, সক্ষমতা ও রক্ষণাবেক্ষণ নিজের মতো পরিচালনা করে। সক্ষমতা বাড়াতে নতুন যন্ত্র ও সময় প্রয়োজন। ক্লাউডে সম্পদ দ্রুত বরাদ্দ করা গেলেও বিল, সংযোগ, সেবাদাতার শর্ত এবং ডেটা রাখার নীতি বিবেচ্য। Pay-as-you-go সবসময় কম মোট খরচ নিশ্চিত করে না। ব্যবহার ও বিল মাপতে হয় এবং অপ্রয়োজনীয় সম্পদ বন্ধ রাখতে হয়।

\subsection{Security ও Shared Responsibility}
\label{ch2:c10-13}
সেবাদাতা ক্লাউড অবকাঠামোর নিরাপত্তা ও পরিচালনার দায়িত্ব নেয়। ব্যবহারকারীর অ্যাকাউন্ট, ডেটা ও প্রবেশাধিকারের দায়িত্ব থাকে। মডেল অনুযায়ী OS ও অ্যাপ্লিকেশন পরিচালনার দায়িত্ব অংশ \ref{ch2:c10-03}-এ আছে। সেবাদাতা আছে বলেই ব্যবহারকারীর সব দায়িত্ব শেষ হয় না।
\chTwoPair{\textbf{কীভাবে সুরক্ষা বাড়তে পারে?}\par
সেবাদাতার প্রবেশাধিকার নিয়ন্ত্রণ, নিয়মিত অবকাঠামো হালনাগাদ ও বিশেষজ্ঞ তদারকি তথ্য রক্ষায় সাহায্য করে। ব্যবহারকারীর শক্তিশালী পাসওয়ার্ড, Multi-factor Authentication (MFA), সীমিত প্রবেশাধিকার, সতর্ক ভাগাভাগি ও উপযুক্ত এনক্রিপশনও প্রয়োজন। MFA-তে ভিন্ন উপাদান দিয়ে অ্যাকাউন্ট যাচাই করা হয়। যেমন, পাসওয়ার্ডের সঙ্গে নিজের ডিভাইসের অনুমোদন।}{\textbf{ঝুঁকি কোথায় থাকে?}\par
অ্যাকাউন্ট চুরি, উন্মুক্ত শেয়ার লিংক, ভুল অনুমতি বা সেবাদাতার ত্রুটিতে অননুমোদিত প্রবেশ ঘটতে পারে। অন্যের সেবায় সংবেদনশীল ডেটা রাখার আগে সংরক্ষণ ও ভাগাভাগির নীতি এবং নিয়ন্ত্রণ দেখতে হয়। অনিরাপদ নেটওয়ার্কে অ্যাকাউন্ট ব্যবহারে সতর্কতা প্রয়োজন। এনক্রিপশন হারানো ফাইল ফেরানোর ব্যবস্থা নয়। ব্যাকআপ ও পুনরুদ্ধারের আলোচনা অংশ \ref{ch2:c10-11}-এ আছে।}

\subsection{Cloud-এর ধারণা, ইতিহাস ও ব্যবহার মাপা}
\chTwoPair{\textbf{ধারণা ও সেবার বিকাশ}\par
John McCarthy ১৯৬১ সালে কম্পিউটিংকে ইউটিলিটি সেবা হিসেবে সংগঠিত করার সম্ভাবনা আলোচনা করেন। পরে নেটওয়ার্ক, ভার্চুয়ালাইজেশন ও ডেটা সেন্টারের উন্নয়নে আধুনিক ক্লাউড গড়ে ওঠে। ২০০৬ সালে Amazon EC2-এর পরীক্ষামূলক Beta সেবা গুরুত্বপূর্ণ বাণিজ্যিক ধাপ। McCarthy-র ধারণা, EC2 চালু ও NIST-এর সংজ্ঞা পৃথক ঘটনা। এগুলো পুরো ক্লাউডের একক যন্ত্র আবিষ্কারের সাল নয়।}{\textbf{একটি কাল্পনিক খরচের হিসাব}\par
স্টোরেজের মাসিক খরচ প্রতি GB ০.৫০ টাকা এবং প্রক্রিয়াকরণের খরচ প্রতি ঘণ্টা ২ টাকা ধরা হলো। ১০০ GB স্টোরেজ ও ১২০ Compute-hour-এর খরচ $100\times0.50+120\times2=290$ টাকা। এটি শেখানোর কাল্পনিক উদাহরণ; সেবাদাতার বর্তমান মূল্য নয়। ডেটা স্থানান্তর, কর ও অন্যান্য খরচ থাকলে পৃথকভাবে যোগ করতে হবে।}
ক্লাউডের জন্য নির্দিষ্ট সর্বোচ্চ দূরত্ব নেই। ক্লায়েন্টের নেটওয়ার্ক সংযোগ ও সেবা পাওয়ার সুযোগ প্রয়োজন। GB সংরক্ষণের পরিমাণ, Mbps ডেটা স্থানান্তরের হার এবং Compute-hour প্রক্রিয়াকরণের ব্যবহৃত সময় বোঝায়। তিনটির একক ও বিল আলাদা। অ্যাকাউন্টসংখ্যা বা বিনামূল্যের স্টোরেজের সীমা সেবাদাতার শর্ত অনুযায়ী বদলায়। একটি সংখ্যাকে ক্লাউডের সংজ্ঞার অংশ ধরা যায় না।
